\documentclass[11pt]{article}
\usepackage[T1]{fontenc}
\usepackage[margin=1in]{geometry}
\usepackage{amsmath,amssymb}
\usepackage[hidelinks]{hyperref}
\usepackage{microtype}

\newcommand{\SL}{\operatorname{SL}}
\newcommand{\Z}{\mathbb Z}
\newcommand{\C}{\mathbb C}

\title{The 2-Kazhdan question for integral special linear groups}
\author{Literature-implied answer note for arXiv:1711.10238}
\date{11 August 2026}

\begin{document}
\maketitle

\paragraph{Status.}
\textbf{Literature-implied answer, complete for \(n\ge3\).}
The large-\(n\) clause is affirmative: \(\SL_n(\Z)\) is 2-Kazhdan for every
\(n\ge4\). The case \(n=3\) is negative.

\section*{Source question}
Question~4.6 of De Chiffre--Glebsky--Lubotzky--Thom
\cite{source}, on source PDF page~24, asks:
\begin{quote}
Is \(\SL_n(\Z)\) 2-Kazhdan (at least for large \(n\))?
\end{quote}
In the source terminology this means
\[
 H^2(\SL_n(\Z),V)=0
 \quad\text{for every unitary representation }V.
\]

\section*{Identification}
Bader--Sauer's Theorem~A \cite{bs}, on supporting PDF page~2, states that
\(\SL_n(\Z)\), as a lattice in \(\SL_n(\mathbb R)\), has property
\((T_{n-2})\). Thus, when \(n\ge4\),
\[
 H^2(\SL_n(\Z),W)=0
 \quad\text{whenever }W^{\SL_n(\Z)}=0.
\]
Their displayed formula~(1) on the same page shows
\(H_c^2(\SL_n(\mathbb R),\C)=0\). Their Theorem~C transfers continuous
cohomology to the lattice in degrees strictly below the real rank \(n-1\),
so
\[
 H^2(\SL_n(\Z),\C)=0 \qquad(n\ge4).
\]

Now split an arbitrary unitary representation orthogonally as
\(V=V^\Gamma\oplus(V^\Gamma)^\perp\), where \(\Gamma=\SL_n(\Z)\).
The second summand has zero degree-two cohomology by property
\((T_{n-2})\). The action on \(V^\Gamma\) is trivial; because \(\Gamma\) is
of finite type in degree two, its finite cochain complex gives
\[
 H^2(\Gamma,V^\Gamma)
 \cong H^2(\Gamma,\C)\otimes V^\Gamma=0.
\]
Hence \(\SL_n(\Z)\) is 2-Kazhdan for every \(n\ge4\).

The endpoint below this range is determined by
Br\"uck--Hughes--Kielak--Mizerka's Theorem~1.1 \cite{bhkm}, on supporting
PDF page~2. They construct a finite-dimensional orthogonal representation
\(\pi_3\) such that
\[
 H^2(\SL_3(\Z),\pi_3)\ne0.
\]
Therefore \(\SL_3(\Z)\) is not 2-Kazhdan.

\section*{Conclusion, provenance, and scope}
Combining the two later papers gives the sharp classification for \(n\ge3\):
\[
 \boxed{\ \SL_3(\Z)\text{ is not 2-Kazhdan, while }
 \SL_n(\Z)\text{ is 2-Kazhdan for every }n\ge4.\ }
\]
This is an agent-identified implication of known theorems, not new
mathematics. The supporting papers work in the source's higher-property-T
framework, but neither presents this exact classification as an answer to
Question~4.6. Exact-question, exact-title, and core-keyword searches through
11 August 2026 led to these decisive sources.

This note does not address the source's separate questions about amenable or
solvable Frobenius approximability, central quotients, semidirect products by
\(\Z\), or cohomology with arbitrary \(C^*\)-algebra coefficients.

\begin{thebibliography}{9}
\raggedright
\bibitem{source}
M. De Chiffre, L. Glebsky, A. Lubotzky, and A. Thom,
\emph{Stability, cohomology vanishing, and non-approximable groups},
Forum Math. Sigma 8 (2020), e18; arXiv:1711.10238.

\bibitem{bs}
U. Bader and R. Sauer,
\emph{Higher Kazhdan property and unitary cohomology of arithmetic groups},
arXiv:2308.06517v3 (2026). See Theorem~A, Theorem~C, and formula~(1).

\bibitem{bhkm}
B. Br\"uck, S. Hughes, D. Kielak, and P. Mizerka,
\emph{Non-vanishing unitary cohomology of low-rank integral special linear
groups}, Int. Math. Res. Not. 2025, no.~15, rnaf230;
arXiv:2410.22310. See Theorem~1.1.
\end{thebibliography}

\end{document}
